% =============================================================================
% SECTION 4 — EXPERIMENTS
% =============================================================================
% Drop-in replacement for the Experiments section of paper_skeleton.tex.
% Uses standard packages already in skeleton: amsmath, booktabs, multirow,
% xcolor, graphicx.
%
% NOTE: Numbers requiring std bands and the unconfirmed 32B / Qwen3-8B
% entries are placed inside \textcolor{red}{...} as visual placeholders.
% =============================================================================

\section{Experiments}
\label{sec:exp}

\subsection{Setup}
\label{sec:exp:setup}

\paragraph{Benchmarks.}
We organise our evaluation around three groups of benchmarks chosen to test
both in-domain agent capability and capability preservation outside the
training domain. The first group, \emph{coding agent benchmarks}, comprises
SWE-Bench-Verified (500 instances) and SWE-Bench-Lite
(300 instances)~\citep{swebench}; both are in-domain for the
post-training pipelines we use and constitute our primary measure of
end-task progress. The second group, \emph{non-agent coding benchmarks},
covers LiveCodeBench~\citep{livecodebench}, OJBench~\citep{ojbench},
and FullStackBench-EN~\citep{fullstackbench}; these probe pure code
generation without an agent loop and serve as a regression check for
capability erosion. The third group, \emph{tool-use and out-of-domain agent
benchmarks}, includes Terminal-Bench~\citep{terminalbench},
$\tau$-bench~\citep{taubench}, and BFCL~\citep{bfcl}; none of
these target software engineering, and the latter two contain no Python
code-editing trajectories at all. They are the appropriate stage on which
to test our central hypothesis that the function-call inductive bias
transfers across task families.

\paragraph{Evaluation protocol.}
All numerical results are obtained on the \emph{final} checkpoint of the
training pipeline. A ``baseline'' row therefore refers to the base model
followed by an agentic post-training pipeline (e.g., r2e-gym), while ``ours''
refers to the same base followed by FIM mid-training and then by the
identical post-training pipeline. We do not directly evaluate
mid-training-only checkpoints: mid-training on FIM data without subsequent
instruction-style post-training degrades instruction-following ability,
making any comparison against instruction-tuned baselines unfair on
benchmarks that require instruction following. Consequently, every gain we
attribute to mid-training should be read as a gain that survives subsequent
post-training. To match each baseline's evaluation setup, the agent
harness is fixed by the post-training pipeline rather than by our method:
runs whose post-training is r2e-gym are evaluated with the r2e-gym scaffold
(an OpenHands fork)~\citep{openhands, r2egym}; runs whose
post-training is swe-smith use SWE-agent~\citep{sweagent}; the Qwen3-8B
runs use OpenHands directly, because the swe-lego post-training data we use
exceeds the 32K context length of r2e-gym/swe-smith scaffolds while
remaining within the 40{,}960-token context of Qwen3.
All percentages are reported as the mean over three independent seeds; the
``\%'' symbol is omitted inside table cells. Maximum turn limits, decoding
temperature and top-$p$, and per-step timeouts are deferred to
Appendix~\ref{app:hp}.

\paragraph{Models and baselines.}
For Qwen2.5-Coder-Instruct at 7B / 14B / 32B~\citep{qwencoder} we
report (i) the instruction-tuned base without any agentic training,
(ii) post-training with r2e-gym reproduced under our setup, (iii) the
official r2e-gym numbers (in grey, where available), (iv) our recipe of
FIM mid-training followed by r2e-gym post-training, and at 7B additionally
the analogous swe-smith~\citep{swesmith} variants. For Qwen3-8B we
follow the swe-lego~\citep{swelego} pipeline; the cross-base-model
comparison thus also varies the post-training pipeline, a confound we
discuss in Section~\ref{sec:exp:main}.


% =============================================================================
% 4.2 Main Results
% =============================================================================
\subsection{Main Results on Agent Benchmarks}
\label{sec:exp:main}

% -----------------------------------------------------------------------------
% TABLE 3
% -----------------------------------------------------------------------------
\begin{table}[t]
\centering
\caption{Main results on coding agent benchmarks (SWE-Bench-Verified and
SWE-Bench-Lite). All cells are percentages. Bolded rows are our recipe
(FIM mid-training followed by the same post-training pipeline as the row
above). Greyed rows are reproduced from the corresponding officially
released numbers; black rows are reproduced under our setup over three
seeds. Std bands shown in red are placeholders awaiting final replication
results; the 32B mid-training entries (red) are still being verified
against intermediate logs.}
\label{tab:main}
\small
\setlength{\tabcolsep}{6pt}
\begin{tabular}{llcc}
\toprule
Base model & Setting & SWE-Bench-Verified & SWE-Bench-Lite \\
\midrule
\multirow{7}{*}{Qwen2.5-Coder-7B-Instruct}
  & --- (no agentic training)            & 1.8\,($\pm$1.3)                              & 1.0\,($\pm$1.0) \\
  & + r2e-gym (reproduced)               & 15.00\,(\textcolor{red}{$\pm$?})              & 11.33\,(\textcolor{red}{$\pm$?}) \\
  & + r2e-gym (official)                 & \textcolor{gray}{19.00}                        & \textcolor{gray}{11.00} \\
  & \textbf{+ FIM-Midtrain + r2e-gym}    & \textbf{17.80}\,(\textcolor{red}{$\pm$?})     & \textbf{15.00}\,(\textcolor{red}{$\pm$?}) \\
\cmidrule(lr){2-4}
  & + swe-smith (reproduced)             & 12.30\,(\textcolor{red}{$\pm$?})              & 14.20\,(\textcolor{red}{$\pm$?}) \\
  & + swe-smith (official)               & \textcolor{gray}{15.20}                        & \textcolor{gray}{11.70} \\
  & \textbf{+ FIM-Midtrain + swe-smith}  & \textbf{17.60}\,(\textcolor{red}{$\pm$?})     & \textbf{14.70}\,(\textcolor{red}{$\pm$?}) \\
\midrule
\multirow{4}{*}{Qwen2.5-Coder-14B-Instruct}
  & --- (no agentic training)            & 4.0\,($\pm$1.6)                              & 2.7\,($\pm$1.0) \\
  & + r2e-gym (reproduced)               & 26.20\,(\textcolor{red}{$\pm$?})              & 18.00\,(\textcolor{red}{$\pm$?}) \\
  & + r2e-gym (official)                 & \textcolor{gray}{26.80}                        & \textcolor{gray}{20.67} \\
  & \textbf{+ FIM-Midtrain + r2e-gym}    & \textbf{29.20}\,(\textcolor{red}{$\pm$?})     & \textbf{22.00}\,(\textcolor{red}{$\pm$?}) \\
\midrule
\multirow{4}{*}{Qwen2.5-Coder-32B-Instruct}
  & --- (no agentic training)            & 7.0\,($\pm$1.3)                              & 3.0\,($\pm$1.4) \\
  & + r2e-gym (reproduced)               & 31.80\,(\textcolor{red}{$\pm$?})              & 24.70\,(\textcolor{red}{$\pm$?}) \\
  & + r2e-gym (official)                 & \textcolor{gray}{34.40}                        & \textcolor{gray}{23.77} \\
  & \textbf{+ FIM-Midtrain + r2e-gym}    & \textbf{\textcolor{red}{35.10}}\,(\textcolor{red}{$\pm$?}) & \textbf{\textcolor{red}{26.80}}\,(\textcolor{red}{$\pm$?}) \\
\midrule
\multirow{2}{*}{Qwen3-8B-Base}
  & + swe-lego                           & 30.20\,(\textcolor{red}{$\pm$?})              & \textcolor{red}{TBD}\,(\textcolor{red}{$\pm$?}) \\
  & \textbf{+ FIM-Midtrain + swe-lego}   & \textbf{\textcolor{red}{33.40}}\,(\textcolor{red}{$\pm$?}) & \textbf{\textcolor{red}{TBD}}\,(\textcolor{red}{$\pm$?}) \\
\bottomrule
\end{tabular}
\end{table}

Table~\ref{tab:main} summarises the in-domain agent results.

\paragraph{Mid-training delivers consistent gains across model scales.}
Holding the post-training pipeline fixed at r2e-gym, FIM mid-training
improves SWE-Bench-Verified by $+2.8$ at 7B ($15.00\!\to\!17.80$),
$+3.0$ at 14B ($26.20\!\to\!29.20$), and $+3.3$ at 32B
($31.80\!\to\!\textcolor{red}{35.10}$). The absolute gain stays within a
narrow band of $+2.8$ to $+3.3$ as the base model grows by more than
$4{\times}$ in parameters. SWE-Bench-Lite shows the same direction at every
scale: $+3.7$ / $+4.0$ / $\textcolor{red}{+2.1}$ at 7B / 14B / 32B. We do
not interpret this as evidence of monotone scaling, but rather as evidence
that the structural prior is not absorbed by larger pretrained models in
the practical deployment range we evaluate. The relative gain naturally
contracts with size ($+18.7\%$ at 7B, $+11.5\%$ at 14B, $+10.4\%$ at 32B),
consistent with the higher absolute baseline at larger scales. We do not
extrapolate these trends beyond 32B.

\paragraph{The improvement transfers across post-training pipelines.}
Replacing r2e-gym with swe-smith on the same 7B base, FIM mid-training
improves SWE-Bench-Verified from $12.30$ to $17.60$, a gain of $+5.3$
points larger than that observed under r2e-gym ($+2.8$). On SWE-Bench-Lite
the swe-smith pairing improves from $14.20$ to $14.70$, a small change of
$+0.5$; we report this honestly because the corresponding r2e-gym pairing
gains $+3.7$ on the same benchmark, so the cross-pipeline transfer is
clearer on Verified than on Lite. Taken together, the two pipelines
indicate that mid-training is not specifically tuned to a single
post-training data distribution, but the magnitude of its benefit can
depend on which baseline pipeline it is composed with.

\paragraph{The improvement transfers across base-model families.}
Switching to Qwen3-8B and the swe-lego post-training pipeline,
mid-training improves SWE-Bench-Verified from $30.20$ to
$\textcolor{red}{33.40}$, a gain of $\textcolor{red}{+3.2}$ points
comparable to those observed on Qwen2.5-Coder. We note that this
comparison varies the post-training pipeline simultaneously with the base
model (swe-lego is required because Qwen3-8B's $40{,}960$-token context
exceeds that of the r2e-gym/swe-smith scaffolds), so the cross-base-model
result should be read as ``the recipe is not specific to the
Qwen2.5-Coder family under one alternative agentic pipeline,'' rather
than as a guarantee that it generalises across all base-model families.
We elaborate this caveat in Section~\ref{sec:limitations}.


% =============================================================================
% 4.3 Capability preservation and cross-domain transfer
% =============================================================================
\subsection{Capability Preservation and Cross-Domain Transfer}
\label{sec:exp:gen}

A natural concern with any post-training stage that specialises a model
toward a narrow agent task is that it may erode capabilities the base
model already possessed. We test for this directly on six benchmarks that
sit outside the SWE-Bench distribution and, in the case of $\tau$-bench
and BFCL, outside the coding domain entirely. Because this analysis
requires a controlled comparison among (instruct base) vs.\ (post-training
only) vs.\ (mid-training $+$ post-training), and the corresponding
controlled triple is the most expensive to produce, we conduct it at the
14B scale with the r2e-gym pipeline. Results are summarised in
Table~\ref{tab:cap}.

% -----------------------------------------------------------------------------
% TABLE 4
% -----------------------------------------------------------------------------
\begin{table}[t]
\centering
\caption{Capability preservation and cross-domain transfer at 14B with the
r2e-gym post-training pipeline. ``Non-agent coding'' tests pure code
generation; ``Agent OOD'' is closer to coding agents but in a different
environment; ``Tool use'' tests non-coding tool-call behaviour. Bold marks
the better trained variant per column (post-only vs.\ ours); the Instruct
row is shown as a reference ceiling and is not in competition. All cells
are percentages, std bands in red are placeholders.}
\label{tab:cap}
\small
\setlength{\tabcolsep}{4pt}
\begin{tabular}{lcccccc|c}
\toprule
& \multicolumn{3}{c}{Non-agent coding}
& \multicolumn{1}{c}{Agent OOD}
& \multicolumn{2}{c}{Tool use}
& \\
\cmidrule(lr){2-4} \cmidrule(lr){5-5} \cmidrule(lr){6-7}
Setting (Qwen2.5-Coder-14B-base)
& LiveCode & OJBench & FSB-EN
& Terminal
& $\tau$-bench & BFCL
& Avg \\
\midrule
Instruct (no agentic training)
  & 37.20 & 5.20 & 53.80
  & 0.00
  & 5.70 & 23.20
  & 20.85 \\
+ r2e-gym
  & 24.10\,(\textcolor{red}{$\pm$?})
  & 2.80\,(\textcolor{red}{$\pm$?})
  & 47.72\,(\textcolor{red}{$\pm$?})
  & 1.28\,(\textcolor{red}{$\pm$?})
  & 0.40\,(\textcolor{red}{$\pm$?})
  & 15.70\,(\textcolor{red}{$\pm$?})
  & 15.33 \\
\textbf{+ FIM-Midtrain + r2e-gym}
  & \textbf{35.20}\,(\textcolor{red}{$\pm$?})
  & \textbf{4.74}\,(\textcolor{red}{$\pm$?})
  & \textbf{48.07}\,(\textcolor{red}{$\pm$?})
  & \textbf{2.47}\,(\textcolor{red}{$\pm$?})
  & \textbf{4.40}\,(\textcolor{red}{$\pm$?})
  & \textbf{18.20}\,(\textcolor{red}{$\pm$?})
  & \textbf{18.85} \\
\bottomrule
\end{tabular}
\end{table}

\paragraph{Agentic post-training has a substantial hidden capability cost.}
On the same 14B base, applying r2e-gym alone reduces LiveCodeBench from
$37.20$ to $24.10$ (a drop of $13.10$ points), OJBench from $5.20$ to
$2.80$, and FullStackBench-EN from $53.80$ to $47.72$. The deterioration
is more severe outside the coding domain: $\tau$-bench falls from $5.70$
to $0.40$ and BFCL from $23.20$ to $15.70$. Averaged across the six
benchmarks the post-trained model loses $5.52$ points relative to the
Instruct ceiling. This regression is rarely highlighted in agent papers,
but it is the implicit cost being paid for SWE-Bench gains.

\paragraph{Mid-training largely closes the regression gap.}
Adding FIM mid-training before the same post-training restores
LiveCodeBench to $35.20$ (recovering $11.10$ of the $13.10$-point drop),
OJBench to $4.74$ (within $0.46$ of the Instruct ceiling), and
FullStackBench-EN to $48.07$. The six-benchmark average rises from
$15.33$ to $18.85$, an increase of $+3.52$ over the post-training-only
configuration, while the SWE-Bench-Verified gain on the same base is
preserved (Table~\ref{tab:main}). In other words, mid-training improves
the cost--benefit profile of agentic post-training: the in-domain target
metric improves, and the bulk of the off-distribution capability erosion
is undone in the same training run.

\paragraph{Cross-domain transfer to non-coding tool use.}
The most direct test of our motivating hypothesis lies in the rightmost
columns of Table~\ref{tab:cap}: $\tau$-bench and BFCL contain no Python
code-editing data, and our mid-training corpus contains no tool-use
trajectories. Mid-training nevertheless improves $\tau$-bench from
$0.40$ to $4.40$ ($+4.0$) and BFCL from $15.70$ to $18.20$ ($+2.5$),
with directionally consistent recovery on the also-OOD Terminal-Bench
($1.28\!\to\!2.47$). Because the mid-training data carries no tool-use
signal a priori, the only mechanism by which it can affect these
benchmarks is by installing a structural prior at mid-training time that
survives subsequent post-training. The fact that this prior moves
$\tau$-bench and BFCL in the same direction as the in-domain SWE-Bench
gains is consistent with our hypothesis that the function-call site and
the agent tool-call step share an underlying conditioning structure
(Section~\ref{sec:method:motivation}).


% =============================================================================
% 4.4 Consistency across model sizes
% =============================================================================
\subsection{Consistency Across Model Sizes}
\label{sec:exp:scaling}

\begin{figure}[t]
\centering
\includegraphics[width=0.92\linewidth]{figs/scaling_curve.pdf}
\caption{SWE-Bench-Verified (left) and SWE-Bench-Lite (right) for the
Qwen2.5-Coder family at 7B, 14B, and 32B, comparing post-training only
(grey) against FIM mid-training followed by the same r2e-gym post-training
pipeline (red). Annotations report the absolute gain from mid-training at
each scale.}
\label{fig:scaling}
\end{figure}

\paragraph{The gain does not vanish with scale.}
Figure~\ref{fig:scaling} replots the in-domain numbers from
Table~\ref{tab:main}. A natural concern with any structural prior is that
larger pretrained models, having seen more code tokens, would already have
absorbed it; in that case mid-training would help only at small scale and
diminish toward $32$B. The data we have do not show this pattern: the
absolute SWE-Bench-Verified gain is $+2.8$ at $7$B, $+3.0$ at $14$B and
$+3.3$ at $32$B, a narrow band of variation as the parameter count grows
$4{\times}$. The same direction holds on SWE-Bench-Lite ($+3.7/+4.0/+2.1$).
The relative gain naturally contracts with size ($+18.7\%$ at $7$B,
$+11.5\%$ at $14$B, $+10.4\%$ at $32$B), consistent with the higher
absolute baseline at larger scales. We do not extrapolate beyond the three
sizes evaluated; whether the trend continues at $70$B+ is left to future
work.


% =============================================================================
% 4.5 Ablation studies
% =============================================================================
\subsection{Ablation Studies}
\label{sec:exp:ablation}

We conduct two controlled ablations, both on the Qwen2.5-Coder-7B base
with r2e-gym post-training and a shared mid-training-free baseline. The
first isolates the role of the chain-of-thought rationale embedded in the
FIM middle span; the second isolates the role of the function-aware
selection pipeline. Due to compute constraints, the ablations run at the
7B scale only; we expect the qualitative ordering to carry to larger
scales but do not verify it. Mask granularity, score-threshold
sensitivity, and mid-training token budget are reserved for future
investigation.

% -----------------------------------------------------------------------------
% TABLE 6  (Ablations)
% -----------------------------------------------------------------------------
\begin{table}[t]
\centering
\caption{Ablation studies on the Qwen2.5-Coder-7B base with r2e-gym
post-training. \textbf{(A)} varies the source of the rationale embedded in
the FIM middle span. \textbf{(B)} varies the function-selection algorithm,
with all other components fixed (CoT source held to Gemini-3 throughout).
The two blocks share the same baseline (\emph{w/o mid-train}) and the same
``full'' configuration, included in both for reference. Bold marks our
default recipe. All cells are percentages.}
\label{tab:ablation}
\small
\setlength{\tabcolsep}{8pt}
\begin{tabular}{lccc}
\toprule
Setting & SWE-Bench-Verified & SWE-Bench-Lite & Average \\
\midrule
\multicolumn{4}{l}{\emph{(A) Chain-of-thought source (selection fixed to ``full'')}} \\
\midrule
w/o mid-train (baseline)            & 15.00 & 11.33 & 13.17 \\
+ FIM, no CoT                       & 15.60 & 11.00 & 13.30 \\
+ FIM, self-CoT (Qwen2.5-Coder-7B)  & 16.40 & 13.30 & 14.85 \\
\textbf{+ FIM, Gemini-3 CoT (full)} & \textbf{17.80} & \textbf{15.00} & \textbf{16.40} \\
\midrule
\multicolumn{4}{l}{\emph{(B) Function-selection algorithm (CoT source fixed to Gemini-3)}} \\
\midrule
w/o mid-train (baseline)            & 15.00 & 11.33 & 13.17 \\
random                              & 15.30 & 12.60 & 13.95 \\
Gemini-selected                     & 16.40 & 13.70 & 15.05 \\
PDG-only                            & 17.00 & 14.70 & 15.85 \\
inferability-only                   & 16.80 & 12.00 & 14.40 \\
\textbf{full (PDG + inferability)}  & \textbf{17.80} & \textbf{15.00} & \textbf{16.40} \\
\bottomrule
\end{tabular}
\end{table}

\paragraph{FIM structure contributes independently of CoT distillation.}
Block (A) of Table~\ref{tab:ablation} disarms a natural reviewer concern:
that our improvements stem primarily from distilling reasoning from a
frontier model. Removing the rationale entirely (\emph{no CoT}) still
improves the average by $+0.13$ over the no-mid-training baseline, with a
$+0.6$ gain on SWE-Bench-Verified, indicating that the function-aware FIM
objective by itself contributes some signal even without any external
rationale. Replacing Gemini-3 with rationales generated by the model
being trained (\emph{self-CoT}) reaches an average of $14.85$, recovering
$1.7$ of the $3.2$-point gap between \emph{no CoT} and \emph{Gemini-3 CoT};
high-quality external rationales help, but a sizeable fraction of the
gain is reproducible without a frontier teacher. We accordingly do not
view the recipe as fundamentally distillation-driven, although our best
configuration does benefit from teacher rationales.

\paragraph{Function-aware selection is the dominant lever within mid-training.}
Block (B) varies the selection algorithm with the rationale source held
fixed to Gemini-3. \emph{Random} masking of functions yields an average of
$13.95$, a margin of only $+0.78$ over no mid-training: most of the
mid-training time spent on randomly chosen functions is wasted. Letting
Gemini-3 directly pick which functions to mask (\emph{Gemini-selected})
reaches $15.05$, comparable to the model-internal \emph{self-CoT}
configuration but markedly worse than our full pipeline. Among the
structural signals, \emph{PDG-only} ($15.85$) outperforms
\emph{inferability-only} ($14.40$), and the two are complementary:
combining them (the \emph{full} configuration) yields $16.40$, a further
$+0.55$ over the better single signal. This is consistent with our design
intent: the program dependency graph identifies functions whose
reconstruction requires non-local reasoning, while the inferability score
ensures the target is recoverable from context---masking targets that
score highly on only one of the two are either unlearnable noise (high
complexity, low inferability) or trivial fills (low complexity, high
inferability).

% -----------------------------------------------------------------------------
% TOKEN-BUDGET ABLATION (placeholder values pending experiment)
% -----------------------------------------------------------------------------
\begin{figure}[t]
\centering
\includegraphics[width=0.55\linewidth]{figs/token_budget.pdf}
\caption{Mid-training token-budget ablation on the Qwen2.5-Coder-14B base
with the r2e-gym pipeline. \textcolor{red}{Numerical values are placeholders
pending the final ablation run.} The dashed line marks the
post-training-only baseline; the circled point is our default $3$B-token
budget. The curve saturates between $3$B and $5$B tokens, suggesting that
additional FIM mid-training data yields diminishing returns within this
range.}
\label{fig:token_budget}
\end{figure}

\paragraph{Token-budget scaling.}
Figure~\ref{fig:token_budget} sweeps the mid-training token budget across
$\{0.5, 1, 2, 3, 5\}$~B tokens at fixed selection algorithm and CoT source.
\textcolor{red}{Pending final numbers, the placeholder curve already
captures the qualitative pattern we expect: returns grow rapidly from
$0.5$~B to $2$~B, the default $3$~B sits near the knee, and the
$5$~B point shows minimal further improvement.} We accordingly adopt
$3$~B tokens as the default in all other experiments to balance gain
against compute.


